\documentclass[10pt]{article}
\usepackage{fontspec}
\IfFontExistsTF{Times New Roman}{\setmainfont{Times New Roman}}{\setmainfont{Liberation Serif}}
\usepackage[russian]{babel}
\usepackage{geometry}
\geometry{a4paper, margin=1.5cm}

\usepackage{longtable, array, booktabs, tabularx}
\usepackage{ragged2e}
\usepackage{fancyhdr}
\usepackage{lastpage}
\usepackage{graphicx}
\usepackage{amsmath}
\usepackage{xcolor}
\usepackage{colortbl}

\newcolumntype{P}[1]{>{\RaggedRight\arraybackslash\footnotesize}m{#1}}
\newcolumntype{C}[1]{>{\centering\arraybackslash\footnotesize}m{#1}}
\newcolumntype{R}[1]{>{\RaggedLeft\arraybackslash\footnotesize}m{#1}}

\setlength{\parindent}{0pt}
\setlength{\parskip}{1pt}
\renewcommand{\arraystretch}{1.15}

\usepackage{titlesec}
\titlespacing*{\section}{0pt}{0.5em}{0.2em}
\titlespacing*{\subsection}{0pt}{0.4em}{0.15em}

\pagestyle{fancy}
\fancyhf{}
\fancyhead[L]{\footnotesize Отчёт об адаптации Скв №\VAR{well.number} · №\VAR{doc_number}}
\fancyhead[R]{\footnotesize \thepage\ / \pageref{LastPage}}
\fancyfoot[C]{\scriptsize Сформировано \VAR{generated_at}}

\begin{document}

% ═══════════════════════════════════════════════════════════════
%  ТИТУЛЬНЫЙ БЛОК
% ═══════════════════════════════════════════════════════════════
\thispagestyle{fancy}

{\large\bfseries Отчёт об адаптации работы скважины №\VAR{well.number}}

\vspace{0.3em}

{\footnotesize
\begin{tabularx}{\textwidth}{XXXX}
\toprule
\textbf{Документ} & \textbf{Горизонт} & \textbf{Штуцер} & \textbf{Сформировано} \\
№\VAR{doc_number} & \VAR{well.horizon} & \VAR{well.choke_mm}\,мм & \VAR{generated_at} \\
\midrule
\textbf{Этап наблюдения} & \textbf{Этап адаптации} & \multicolumn{2}{l}{\textbf{Длительность}} \\
\VAR{observation.date_from}--\VAR{observation.date_to} & \VAR{adaptation.date_from}--\VAR{adaptation.date_to} & \multicolumn{2}{l}{Набл. \VAR{observation.duration_days} сут., адапт. \VAR{adaptation.duration_days} сут.} \\
\bottomrule
\end{tabularx}
}

\BLOCK{ if warnings }
\vspace{0.3em}
{\scriptsize\color{orange!70!black}
\textbf{Предупреждения:}
\begin{itemize}\setlength{\itemsep}{-2pt}
%% for w in warnings
\item \VAR{w}
%% endfor
\end{itemize}}
\BLOCK{ endif }

\vspace{0.5em}

% ═══════════════════════════════════════════════════════════════
%  1. ТЕХНИЧЕСКИЕ ДАННЫЕ  (toggle: include_sections.well_info)
% ═══════════════════════════════════════════════════════════════
\BLOCK{ if include_sections.well_info }
\section*{1. Технические данные}

% ── 1.1. Сводная таблица параметров ──
{\scriptsize Табл.~1.1. Сводные параметры.}
\vspace{0.2em}

{\footnotesize
\begin{tabularx}{\textwidth}{|p{0.30\textwidth}|X|}
\hline
\textbf{Параметр} & \textbf{Значение} \\
\hline
Номер скважины & \VAR{well.number} \\
\hline
Наименование & \VAR{well.name} \\
\hline
Пласт / Горизонт & \VAR{well.horizon} \\
\hline
Диаметр штуцера, мм & \VAR{well.choke_mm}
\BLOCK{ if well.choke_history and well.choke_history|length > 1 }
\newline
{\scriptsize\textit{история:
\BLOCK{ for h in well.choke_history }\VAR{h.data_as_of or '—'}~{:}~\VAR{'%g'|format(h.choke_mm)}\,мм\BLOCK{ if not loop.last }; \BLOCK{ endif }\BLOCK{ endfor }}}
\BLOCK{ endif } \\
\hline
Дата начала наблюдения & \VAR{observation.date_from} \\
\hline
\end{tabularx}
}

% ── Карта на всю ширину, высотой не более 30% страницы ──
\BLOCK{ if well.map_chart_path }
\vspace{0.4em}

\begin{center}
\includegraphics[width=\textwidth,height=0.30\textheight,keepaspectratio]{\VAR{well.map_chart_path}}
\end{center}
{\scriptsize\centering Рис.~1.1. Расположение относительно соседних скважин.
\textcolor{red!70!black}{\textbf{★}}~--- скв. №\VAR{well.number};
\textcolor{gray!70!black}{\textbullet}~--- соседние скважины.\par}
\BLOCK{ endif }

\vspace{0.4em}

% ── 1.2. Конструкция скважины ──
\BLOCK{ if well.construction }
\subsection*{1.2. Конструкция скважины}

{\scriptsize Табл.~1.2. Параметры конструкции.}
\vspace{0.2em}

{\footnotesize
\begin{tabularx}{\textwidth}{|p{0.68\textwidth}|X|}
\hline
\textbf{Параметр} & \textbf{Значение} \\
\hline
\BLOCK{ if well.construction.prod_casing_diam_mm is not none }
Диаметр эксплуатационной колонны, мм & \VAR{'%g'|format(well.construction.prod_casing_diam_mm)} \\
\hline
\BLOCK{ endif }
\BLOCK{ if well.construction.prod_casing_depth_m is not none }
Глубина спуска эксплуатационной колонны, м & \VAR{'%g'|format(well.construction.prod_casing_depth_m)} \\
\hline
\BLOCK{ endif }
\BLOCK{ if well.construction.current_bottomhole_m is not none }
Текущий забой, м & \VAR{'%g'|format(well.construction.current_bottomhole_m)} \\
\hline
\BLOCK{ endif }
\BLOCK{ if well.construction.tubing_diam_mm is not none }
Диаметр НКТ, мм & \VAR{'%g'|format(well.construction.tubing_diam_mm)} \\
\hline
\BLOCK{ endif }
\BLOCK{ if well.construction.tubing_shoe_depth_m is not none }
Глубина башмака НКТ, м & \VAR{'%g'|format(well.construction.tubing_shoe_depth_m)} \\
\hline
\BLOCK{ endif }
\BLOCK{ if well.construction.packer_depth_m is not none }
Глубина пакера, м & \VAR{'%g'|format(well.construction.packer_depth_m)} \\
\hline
\BLOCK{ endif }
\BLOCK{ if well.construction.adapter_depth_m is not none }
Глубина переводника, м & \VAR{'%g'|format(well.construction.adapter_depth_m)} \\
\hline
\BLOCK{ endif }
\BLOCK{ if well.construction.pattern_stuck_depth_m is not none }
Глубина непрохода шаблона, м & \VAR{'%g'|format(well.construction.pattern_stuck_depth_m)} \\
\hline
\BLOCK{ endif }
\BLOCK{ if well.construction.data_as_of }
Дата актуализации & \VAR{well.construction.data_as_of} \\
\hline
\BLOCK{ endif }
\end{tabularx}
}

\BLOCK{ if well.perforation and well.perforation|length > 0 }
\vspace{0.3em}
{\scriptsize Табл.~1.3. Интервалы перфорации.}
\vspace{0.2em}

{\footnotesize
\begin{tabularx}{\textwidth}{|C{2cm}|R{4cm}|R{4cm}|X|}
\hline
\textbf{№} & \textbf{Сверху, м} & \textbf{Снизу, м} & \textbf{Мощность, м} \\
\hline
%% for p in well.perforation
\VAR{p.interval_index}
& \BLOCK{ if p.top_depth_m is not none }\VAR{'%g'|format(p.top_depth_m)}\BLOCK{ else }---\BLOCK{ endif }
& \BLOCK{ if p.bottom_depth_m is not none }\VAR{'%g'|format(p.bottom_depth_m)}\BLOCK{ else }---\BLOCK{ endif }
& \BLOCK{ if p.top_depth_m is not none and p.bottom_depth_m is not none }\VAR{'%g'|format(p.bottom_depth_m - p.top_depth_m)}\BLOCK{ else }---\BLOCK{ endif } \\
\hline
%% endfor
\end{tabularx}
}
\BLOCK{ endif }
\BLOCK{ endif }
%# /1.2. Конструкция скважины

\BLOCK{ if equipment_acts }
\vspace{0.3em}
{\scriptsize Табл.~1.4. Установленное оборудование (акты).}
\vspace{0.2em}

{\footnotesize
\begin{tabularx}{\textwidth}{|C{2.5cm}|C{2cm}|X|}
\hline
\textbf{Номер акта} & \textbf{Дата} & \textbf{Оборудование} \\
\hline
%% for act in equipment_acts
\VAR{act.act_number} & \VAR{act.act_date} & \VAR{act.equipment_name} \\
\hline
%% endfor
\end{tabularx}
}
\BLOCK{ endif }

\vspace{0.3em}
\BLOCK{ endif }
%# /1. Технические данные

% ═══════════════════════════════════════════════════════════════
%  2. АНАЛИЗ ИСХОДНЫХ ДАННЫХ  (toggle: include_sections.customer_data)
% ═══════════════════════════════════════════════════════════════
\BLOCK{ if include_sections.customer_data and customer_chapter and customer_chapter.periods }
\section*{2. Анализ исходных данных (по данным заказчика)}

{\footnotesize
Источник --- суточные сводки заказчика (таблица well\_daily, выгрузка УзКорГаз).
Шаг измерений: 1 запись в сутки. Анализируемые параметры:
$P_\text{уст}$, $P_\text{затр}$, $P_\text{лин}$, $P_\text{статич}$ (кгс/см²),
$Q$ общий и рабочий (тыс.м³/сут), простой (мин/сут), штуцер (мм).
$\Delta P = P_\text{уст} - P_\text{лин}$ (по строкам, неотрицательно).
Тренды считаются как наклон линейной регрессии параметра по времени
(единица --- ед.~параметра~/~сут.).
В этой главе данные только описываются; сравнений с baseline нет.
В отчёт включено \VAR{customer_chapter.periods | length} период(а/ов).
}

\vspace{0.3em}

% Каждый период анализа отдельной подсекцией
%% for p in customer_chapter.periods
\subsection*{2.\VAR{loop.index}. Период \VAR{p.period_from}~--~\VAR{p.period_to}}

% Сводка периода (динамические числа)
\BLOCK{ if p.period_summary }
{\footnotesize \VAR{p.period_summary}}

\vspace{0.3em}
\BLOCK{ endif }

\BLOCK{ if p.description_tex }
{\footnotesize \textit{\VAR{p.description_tex}}}

\vspace{0.3em}
\BLOCK{ endif }

% ─── Графики (4 шт.): давления, ΔP, дебиты+простой, помесячно ───
\BLOCK{ if p.chart_paths and p.chart_paths.pressures }
\vspace{0.2em}
\begin{center}
\includegraphics[width=\textwidth,height=0.28\textheight,keepaspectratio]{\VAR{p.chart_paths.pressures}}
\end{center}
{\scriptsize \textbf{Рис.~2.\VAR{loop.index}.1.} Давления (P\textsubscript{уст}, P\textsubscript{затр}, P\textsubscript{лин}, P\textsubscript{статич}) по суткам.\BLOCK{ if p.chart_captions and p.chart_captions.pressures } \VAR{p.chart_captions.pressures}\BLOCK{ endif }}
\BLOCK{ endif }

\BLOCK{ if p.chart_paths and p.chart_paths.dp }
\vspace{0.3em}
\begin{center}
\includegraphics[width=\textwidth,height=0.25\textheight,keepaspectratio]{\VAR{p.chart_paths.dp}}
\end{center}
{\scriptsize \textbf{Рис.~2.\VAR{loop.index}.2.} Перепад давления $\Delta P$.\BLOCK{ if p.chart_captions and p.chart_captions.dp } \VAR{p.chart_captions.dp}\BLOCK{ endif }}
\BLOCK{ endif }

\BLOCK{ if p.chart_paths and p.chart_paths.flow_dt }
\vspace{0.3em}
\begin{center}
\includegraphics[width=\textwidth,height=0.28\textheight,keepaspectratio]{\VAR{p.chart_paths.flow_dt}}
\end{center}
{\scriptsize \textbf{Рис.~2.\VAR{loop.index}.3.} Дебиты (Q общий и рабочий) и простой по суткам.\BLOCK{ if p.chart_captions and p.chart_captions.flow_dt } \VAR{p.chart_captions.flow_dt}\BLOCK{ endif }}
\BLOCK{ endif }

\BLOCK{ if p.chart_paths and p.chart_paths.monthly }
\vspace{0.3em}
\begin{center}
\includegraphics[width=\textwidth,keepaspectratio]{\VAR{p.chart_paths.monthly}}
\end{center}
{\scriptsize \textbf{Рис.~2.\VAR{loop.index}.4.} Помесячно: общий и рабочий дебит (среднее и медиана).}
\BLOCK{ endif }

% ─── Описательная статистика за период ───
\BLOCK{ if p.describe and p.describe|length > 0 }
\vspace{0.5em}
{\scriptsize \textbf{Табл.~2.\VAR{loop.index}.1.} Описательная статистика за период (\VAR{p.days_count or 0}~сут.).}
\vspace{0.2em}

{\scriptsize
\begin{tabularx}{\textwidth}{|X|R{0.9cm}|R{1.6cm}|R{1.4cm}|R{1.4cm}|R{1.6cm}|R{1.4cm}|}
\hline
\textbf{Параметр} & \textbf{N} & \textbf{Среднее} & \textbf{СКО} &
\textbf{Мин} & \textbf{Медиана} & \textbf{Макс} \\
\hline
%% for r in p.describe
\VAR{r.label} & \VAR{r.n}
& \BLOCK{ if r.mean   is not none }\VAR{'%.2f'|format(r.mean)}\BLOCK{ else }---\BLOCK{ endif }
& \BLOCK{ if r.std    is not none }\VAR{'%.2f'|format(r.std)}\BLOCK{ else }---\BLOCK{ endif }
& \BLOCK{ if r.min    is not none }\VAR{'%.2f'|format(r.min)}\BLOCK{ else }---\BLOCK{ endif }
& \BLOCK{ if r.median is not none }\VAR{'%.2f'|format(r.median)}\BLOCK{ else }---\BLOCK{ endif }
& \BLOCK{ if r.max    is not none }\VAR{'%.2f'|format(r.max)}\BLOCK{ else }---\BLOCK{ endif } \\
\hline
%% endfor
\end{tabularx}
}

\vspace{0.1em}
{\scriptsize \textit{Пояснения:} \textbf{N} --- число суточных записей с непустым значением;
\textbf{СКО} --- среднеквадратичное отклонение (разброс относительно среднего).
Медиана устойчивее к выбросам, чем среднее.}
\BLOCK{ endif }

% ─── Помесячная агрегация и тренды ───
\BLOCK{ if p.months and p.months|length > 0 }
\vspace{0.4em}
{\scriptsize \textbf{Табл.~2.\VAR{loop.index}.2.} Помесячная агрегация и тренды.}
\vspace{0.2em}

{\scriptsize
\begin{tabularx}{\textwidth}{|X|R{0.8cm}|R{1.2cm}|R{1.2cm}|R{1.2cm}|R{1.2cm}|R{1.0cm}|R{1.0cm}|R{1.4cm}|R{1.4cm}|}
\hline
\textbf{Месяц} & \textbf{Дней}
& \textbf{Q общ. ср.} & \textbf{Q общ. мед.}
& \textbf{Q раб. ср.} & \textbf{Q раб. мед.}
& \textbf{ΔP ср.} & \textbf{ΔP мед.}
& \textbf{Тренд Q} & \textbf{Тренд ΔP} \\
\hline
%% for m in p.months
\VAR{m.month_label} & \VAR{m.days}
& \BLOCK{ if m.mean_q_total     is not none }\VAR{'%.2f'|format(m.mean_q_total)}\BLOCK{ else }---\BLOCK{ endif }
& \BLOCK{ if m.median_q_total   is not none }\VAR{'%.2f'|format(m.median_q_total)}\BLOCK{ else }---\BLOCK{ endif }
& \BLOCK{ if m.mean_q_working   is not none }\VAR{'%.2f'|format(m.mean_q_working)}\BLOCK{ else }---\BLOCK{ endif }
& \BLOCK{ if m.median_q_working is not none }\VAR{'%.2f'|format(m.median_q_working)}\BLOCK{ else }---\BLOCK{ endif }
& \BLOCK{ if m.mean_dp          is not none }\VAR{'%.2f'|format(m.mean_dp)}\BLOCK{ else }---\BLOCK{ endif }
& \BLOCK{ if m.median_dp        is not none }\VAR{'%.2f'|format(m.median_dp)}\BLOCK{ else }---\BLOCK{ endif }
& \BLOCK{ if m.trend_q_total    is not none }\VAR{'%+.3f'|format(m.trend_q_total)}\BLOCK{ else }---\BLOCK{ endif }
& \BLOCK{ if m.trend_dp         is not none }\VAR{'%+.3f'|format(m.trend_dp)}\BLOCK{ else }---\BLOCK{ endif } \\
\hline
%% endfor
\end{tabularx}
}

\vspace{0.1em}
{\scriptsize \textit{Пояснения:} тренды --- наклон линейной регрессии за месяц
(\textbf{Тренд Q} --- прирост тыс.м³/сут за день; \textbf{Тренд ΔP} --- прирост кгс/см² за день).
Положительное значение --- рост за месяц, отрицательное --- снижение.}
\BLOCK{ endif }

% ─── Описание по месяцам (помесячный анализ) ───
\BLOCK{ if p.month_descriptions and p.month_descriptions|length > 0 }
\vspace{0.3em}
{\footnotesize\textbf{Помесячный анализ}}
\vspace{0.1em}

{\footnotesize
\begin{itemize}\setlength{\itemsep}{2pt}
%% for d in p.month_descriptions
\item \textbf{\VAR{d.label}.} \VAR{d.text}
%% endfor
\end{itemize}
}
\BLOCK{ endif }

\vspace{0.5em}

%% endfor

\newpage

\BLOCK{ else }
% Заглушка если глава не включена
\BLOCK{ endif }

% ═══════════════════════════════════════════════════════════════
%  3. ЭТАП НАБЛЮДЕНИЯ  (toggle: include_sections.observation)
% ═══════════════════════════════════════════════════════════════
\BLOCK{ if include_sections.observation }
\section*{3. Этап наблюдения}

\subsection*{3.1. Период и описание}

\BLOCK{ if observation.intro and observation.intro.intro_text }
{\footnotesize \VAR{observation.intro.intro_text}}

\vspace{0.3em}
\BLOCK{ endif }

\BLOCK{ if observation.intro and observation.intro.sensors and observation.intro.sensors|length > 0 }
{\scriptsize \textbf{Табл.~3.0.} Датчики давления, установленные на скважине.}
\vspace{0.2em}

{\scriptsize
\begin{tabularx}{\textwidth}{|R{3cm}|X|R{2.4cm}|R{2.6cm}|}
\hline
\textbf{Серийный номер} & \textbf{Метка} & \textbf{Положение} & \textbf{Дата установки} \\
\hline
%% for s in observation.intro.sensors
\VAR{s.serial} & \VAR{s.label or '---'} & \VAR{s.position_label} & \VAR{s.installed_at_fmt or '---'} \\
\hline
%% endfor
\end{tabularx}
}

\vspace{0.3em}
\BLOCK{ endif }

{\footnotesize
Период наблюдения: \VAR{observation.date_from}~---~\VAR{observation.date_to} (\VAR{observation.duration_days}~сут.).
\BLOCK{ if observation.choke_label }

Штуцер на этапе: \VAR{observation.choke_label}.
\BLOCK{ endif }
\BLOCK{ if observation.description }

Описание: \VAR{observation.description}
\BLOCK{ endif }
}

\vspace{0.2em}

\subsection*{3.2. Рабочие параметры}

{\scriptsize Табл.~3.1. Рабочие параметры этапа наблюдения.}
\vspace{0.2em}

{\footnotesize
\begin{tabularx}{\textwidth}{|P{4.2cm}|R{2cm}|C{2cm}|X|}
\hline
\textbf{Параметр} & \textbf{Значение} & \textbf{Единица} & \textbf{Примечание} \\
\hline
$P_{\text{тр}}$ (медиана) & \VAR{observation.p_tube_median_fmt} & кгс/см$^2$ & Трубное давление до штуцера, медиана за этап \\
\hline
$P_{\text{лин}}$ (медиана) & \VAR{observation.p_line_median_fmt} & кгс/см$^2$ & Линейное давление после штуцера, медиана за этап \\
\hline
$\Delta P$ (медиана) & \VAR{observation.dp_median_fmt} & кгс/см$^2$ & Перепад $P_{\text{тр}} - P_{\text{лин}}$, медиана за этап \\
\hline
\BLOCK{ if observation.dp_trend }
Тренд $\Delta P$ & \VAR{observation.dp_trend_fmt} & кгс/см$^2$/сут &
\BLOCK{ if observation.dp_trend.direction == "up" }$\uparrow$ Рост перепада: улучшение работы скважины (более стабильный режим, рост дебита)
\BLOCK{ elif observation.dp_trend.direction == "down" }$\downarrow$ Снижение перепада: ослабление работы скважины
\BLOCK{ else }$\rightarrow$ Стабильно, без выраженного тренда\BLOCK{ endif } \\
\hline
\BLOCK{ endif }
$Q$ медианный & \VAR{observation.flow_median_fmt} & тыс.м$^3$/сут & Медиана мгновенного дебита (устойчива к выбросам) \\
\hline
$Q$ средний & \VAR{observation.flow_avg_fmt} & тыс.м$^3$/сут & Среднее арифметическое мгновенного дебита \\
\hline
$\Sigma Q$ накопленный & \VAR{observation.flow_cumulative_fmt} & тыс.м$^3$ & Накопленная добыча за период этапа \\
\hline
$T_{\text{раб}}$ & \VAR{observation.working_hours_fmt} & ч & Часов в работе (с $\Delta P > 0{,}1$ кгс/см$^2$) \\
\hline
КИВ & \VAR{observation.utilization_pct_fmt}\% & --- & Коэфф. использования времени: $T_{\text{раб}}$ / длительность этапа \\
\hline
Простои & \VAR{observation.downtime_hours_fmt} & ч & Суммарное время простоев за этап \\
\hline
$L_{\text{простой}}$ & \VAR{observation.downtime_loss_fmt} & тыс.м$^3$ & Недополученная добыча из-за простоев \\
\hline
\end{tabularx}
}

\vspace{0.3em}

\subsection*{3.3. События и реагенты}

{\scriptsize Табл.~3.2. События за этап наблюдения.}
\vspace{0.2em}

{\footnotesize
\begin{tabularx}{\textwidth}{|X|R{3cm}|R{3cm}|}
\hline
\textbf{Тип} & \textbf{Количество} & \textbf{Расход} \\
\hline
Продувки & \VAR{observation.purge_count} & --- \\
\hline
Вбросы реагентов (всего) & \VAR{observation.reagent_count} & \VAR{observation.reagent_qty_fmt} \\
\hline
\BLOCK{ for r in observation.reagent_stats }
\quad \VAR{r.name} & \VAR{r.count} & \VAR{r.total_qty_fmt} \\
\hline
\BLOCK{ endfor }
\end{tabularx}
}

\vspace{0.4em}

\subsection*{3.4. Анализ этапа наблюдения}

\BLOCK{ if observation.narrative }
{\footnotesize
\VAR{observation.narrative}
}
\BLOCK{ endif }

% ── Графики этапа наблюдения на отдельной странице ──
\BLOCK{ if observation.combined_chart_path }
\newpage
\subsection*{3.5. Графики этапа наблюдения}

\begin{center}
\includegraphics[width=\textwidth,height=0.85\textheight,keepaspectratio]{\VAR{observation.combined_chart_path}}
\end{center}

{\scriptsize \textbf{Рис.~3.1.} Давления (трубное и линейное), перепад $\Delta P$ и мгновенный дебит $Q(t)$ за этап наблюдения.
\BLOCK{ if observation.p_tube_median is not none }Медиана $P_\text{тр}$ \VAR{'%.2f'|format(observation.p_tube_median)}~кгс/см²\BLOCK{ endif }\BLOCK{ if observation.p_line_median is not none }, $P_\text{лин}$ \VAR{'%.2f'|format(observation.p_line_median)}~кгс/см²\BLOCK{ endif }\BLOCK{ if observation.dp_median is not none }, $\Delta P$ \VAR{'%.2f'|format(observation.dp_median)}~кгс/см²\BLOCK{ endif }.
\BLOCK{ if observation.flow_median is not none }Медианный $Q$ \VAR{'%.2f'|format(observation.flow_median)}, средний \VAR{'%.2f'|format(observation.flow_avg)}~тыс.м³/сут\BLOCK{ endif }\BLOCK{ if observation.flow_cumulative is not none }, накопленная добыча \VAR{'%.1f'|format(observation.flow_cumulative)}~тыс.м³ за период\BLOCK{ endif }.
Маркеры: зелёные точки --- вбросы реагента, фиолетовый пунктир --- продувки, серые линии --- прочие события.}
\BLOCK{ endif }

% ── 3.6. Сравнение этапа наблюдения с базовым периодом (если выбран) ──
\BLOCK{ if obs_vs_baseline and obs_vs_baseline.baseline }
\newpage
\subsection*{3.6. Сравнение с базовым периодом (по данным заказчика)}

{\footnotesize
В качестве базового периода (до начала работ по адаптации) выбран:
\textbf{\VAR{obs_vs_baseline.baseline.name}}
(\VAR{obs_vs_baseline.baseline.period_from}~--~\VAR{obs_vs_baseline.baseline.period_to},
\VAR{obs_vs_baseline.baseline.days_count or 0}~сут.).
Сравнение характерных значений этапа наблюдения с базовым периодом.
}

\vspace{0.3em}

{\scriptsize \textbf{Табл.~3.3.} Сравнение этапа наблюдения с базовым периодом.}
\vspace{0.2em}

{\footnotesize
\begin{tabularx}{\textwidth}{|X|R{2.4cm}|R{2.4cm}|R{2.0cm}|R{1.8cm}|}
\hline
\textbf{Параметр} & \textbf{Базовый период} & \textbf{Этап наблюдения} & \textbf{$\Delta$} & \textbf{$\Delta$, \%} \\
\hline
$Q$ медианный, тыс.м³/сут
& \BLOCK{ if obs_vs_baseline.baseline.q_total_median is not none }\VAR{'%.2f'|format(obs_vs_baseline.baseline.q_total_median)}\BLOCK{ else }---\BLOCK{ endif }
& \VAR{observation.flow_median_fmt}
& \BLOCK{ if obs_vs_baseline.customer.delta_q_med is not none }\VAR{'%+.2f'|format(obs_vs_baseline.customer.delta_q_med)}\BLOCK{ else }---\BLOCK{ endif }
& \BLOCK{ if obs_vs_baseline.customer.delta_q_med_pct is not none }\VAR{'%+.1f'|format(obs_vs_baseline.customer.delta_q_med_pct)}\BLOCK{ else }---\BLOCK{ endif } \\
\hline
$\Delta P$ медианный, кгс/см²
& \BLOCK{ if obs_vs_baseline.baseline.dp_median is not none }\VAR{'%.2f'|format(obs_vs_baseline.baseline.dp_median)}\BLOCK{ else }---\BLOCK{ endif }
& \VAR{observation.dp_median_fmt}
& \BLOCK{ if obs_vs_baseline.customer.delta_dp_med is not none }\VAR{'%+.2f'|format(obs_vs_baseline.customer.delta_dp_med)}\BLOCK{ else }---\BLOCK{ endif }
& \BLOCK{ if obs_vs_baseline.customer.delta_dp_med_pct is not none }\VAR{'%+.1f'|format(obs_vs_baseline.customer.delta_dp_med_pct)}\BLOCK{ else }---\BLOCK{ endif } \\
\hline
$P_\text{уст}$ медианный, кгс/см²
& \BLOCK{ if obs_vs_baseline.baseline.p_wellhead_median is not none }\VAR{'%.2f'|format(obs_vs_baseline.baseline.p_wellhead_median)}\BLOCK{ else }---\BLOCK{ endif }
& \VAR{observation.p_tube_median_fmt}
& \BLOCK{ if obs_vs_baseline.customer.delta_p_tube_med is not none }\VAR{'%+.2f'|format(obs_vs_baseline.customer.delta_p_tube_med)}\BLOCK{ else }---\BLOCK{ endif }
& \BLOCK{ if obs_vs_baseline.customer.delta_p_tube_med_pct is not none }\VAR{'%+.1f'|format(obs_vs_baseline.customer.delta_p_tube_med_pct)}\BLOCK{ else }---\BLOCK{ endif } \\
\hline
$P_\text{лин}$ медианный, кгс/см²
& \BLOCK{ if obs_vs_baseline.baseline.p_flowline_median is not none }\VAR{'%.2f'|format(obs_vs_baseline.baseline.p_flowline_median)}\BLOCK{ else }---\BLOCK{ endif }
& \VAR{observation.p_line_median_fmt}
& \BLOCK{ if obs_vs_baseline.customer.delta_p_line_med is not none }\VAR{'%+.2f'|format(obs_vs_baseline.customer.delta_p_line_med)}\BLOCK{ else }---\BLOCK{ endif }
& \BLOCK{ if obs_vs_baseline.customer.delta_p_line_med_pct is not none }\VAR{'%+.1f'|format(obs_vs_baseline.customer.delta_p_line_med_pct)}\BLOCK{ else }---\BLOCK{ endif } \\
\hline
Простой, ч
& \BLOCK{ if obs_vs_baseline.baseline.shutdown_min_total is not none }\VAR{'%.1f'|format(obs_vs_baseline.baseline.shutdown_min_total / 60.0)}\BLOCK{ else }---\BLOCK{ endif }
& \VAR{observation.downtime_hours_fmt}
& \BLOCK{ if obs_vs_baseline.customer.delta_downtime is not none }\VAR{'%+.1f'|format(obs_vs_baseline.customer.delta_downtime)}\BLOCK{ else }---\BLOCK{ endif }
& \BLOCK{ if obs_vs_baseline.customer.delta_downtime_pct is not none }\VAR{'%+.1f'|format(obs_vs_baseline.customer.delta_downtime_pct)}\BLOCK{ else }---\BLOCK{ endif } \\
\hline
\end{tabularx}
}

\vspace{0.2em}
{\scriptsize \textit{Пояснения:} столбец $\Delta$ --- разность «этап наблюдения минус базовый период» по медианным значениям;
$\Delta\,\%$ --- то же относительно базы. Положительный $\Delta Q$ --- скважина даёт больше; положительный $\Delta(\Delta P)$ --- перепад вырос.}

\BLOCK{ endif }
%# /3.6 сравнение с базовым периодом

\newpage
\BLOCK{ endif }
%# /3. Этап наблюдения

% ═══════════════════════════════════════════════════════════════
%  4. ЭТАП АДАПТАЦИИ  (toggle: include_sections.adaptation)
% ═══════════════════════════════════════════════════════════════
\BLOCK{ if include_sections.adaptation }
\section*{4. Этап адаптации}

\subsection*{4.1. Период и описание}

{\footnotesize
Период: \VAR{adaptation.date_from}~---~\VAR{adaptation.date_to} (\VAR{adaptation.duration_days}~сут.).
\BLOCK{ if adaptation.choke_label }

Штуцер на этапе: \VAR{adaptation.choke_label}.
\BLOCK{ endif }
\BLOCK{ if adaptation.description }

Описание: \VAR{adaptation.description}
\BLOCK{ endif }
}

\vspace{0.2em}

\subsection*{4.2. Рабочие параметры}

{\scriptsize Табл.~4.1. Рабочие параметры этапа адаптации.}
\vspace{0.2em}

{\footnotesize
\begin{tabularx}{\textwidth}{|X|R{2.2cm}|R{2.2cm}|R{2.2cm}|}
\hline
\textbf{Параметр} & \textbf{Значение} & \textbf{Единица} & \textbf{Примечание} \\
\hline
$P_{\text{тр}}$ (медиана) & \VAR{adaptation.p_tube_median_fmt} & кгс/см$^2$ & Давление трубное \\
\hline
$P_{\text{лин}}$ (медиана) & \VAR{adaptation.p_line_median_fmt} & кгс/см$^2$ & Давление линейное \\
\hline
$\Delta P$ (медиана) & \VAR{adaptation.dp_median_fmt} & кгс/см$^2$ & Перепад \\
\hline
\BLOCK{ if adaptation.dp_trend }
Тренд $\Delta P$ & \VAR{adaptation.dp_trend_fmt} & кгс/см$^2$/сут &
\BLOCK{ if adaptation.dp_trend.direction == "up" }$\uparrow$ Рост
\BLOCK{ elif adaptation.dp_trend.direction == "down" }$\downarrow$ Снижение
\BLOCK{ else }$\rightarrow$ Стабильно\BLOCK{ endif } \\
\hline
\BLOCK{ endif }
$Q$ медианный & \VAR{adaptation.flow_median_fmt} & тыс.м$^3$/сут & --- \\
\hline
$Q$ средний & \VAR{adaptation.flow_avg_fmt} & тыс.м$^3$/сут & --- \\
\hline
$\Sigma Q$ накопленный & \VAR{adaptation.flow_cumulative_fmt} & тыс.м$^3$ & За период этапа \\
\hline
$T_{\text{раб}}$ & \VAR{adaptation.working_hours_fmt} & ч & Часов с $\Delta P > 0.1$ \\
\hline
КИВ & \VAR{adaptation.utilization_pct_fmt}\% & --- & Коэфф. использ. времени \\
\hline
Простои & \VAR{adaptation.downtime_hours_fmt} & ч & --- \\
\hline
$L_{\text{простой}}$ & \VAR{adaptation.downtime_loss_fmt} & тыс.м$^3$ & Потери от простоев \\
\hline
\end{tabularx}
}

\vspace{0.3em}

\subsection*{4.3. События и реагенты}

{\scriptsize Табл.~4.2. События за этап адаптации.}
\vspace{0.2em}

{\footnotesize
\begin{tabularx}{\textwidth}{|X|R{3cm}|R{3cm}|}
\hline
\textbf{Тип} & \textbf{Количество} & \textbf{Расход} \\
\hline
Продувки & \VAR{adaptation.purge_count} & --- \\
\hline
Вбросы реагентов (всего) & \VAR{adaptation.reagent_count} & \VAR{adaptation.reagent_qty_fmt} \\
\hline
\BLOCK{ for r in adaptation.reagent_stats }
\quad \VAR{r.name} & \VAR{r.count} & \VAR{r.total_qty_fmt} \\
\hline
\BLOCK{ endfor }
\end{tabularx}
}

\vspace{0.4em}

\subsection*{4.4. Анализ этапа адаптации}

\BLOCK{ if adaptation.narrative }
{\footnotesize
\VAR{adaptation.narrative}
}
\BLOCK{ endif }

% ── Графики этапа адаптации на отдельной странице ──
\BLOCK{ if adaptation.combined_chart_path }
\newpage
\subsection*{4.5. Графики этапа адаптации}

\begin{center}
\includegraphics[width=\textwidth,height=0.85\textheight,keepaspectratio]{\VAR{adaptation.combined_chart_path}}
\end{center}

{\scriptsize Рис.~4.1. Давления, перепад $\Delta P$ и мгновенный дебит Q(t) за этап адаптации.
Маркеры: зелёные точки --- вбросы реагента, фиолетовый пунктир --- продувки, серые линии --- прочие события.}
\BLOCK{ endif }

% ── Оптимальный режим (подпериод внутри адаптации) ──
\BLOCK{ if optimal_regime }
\newpage
\subsection*{4.6. Оптимальный режим работы}

{\footnotesize
Оптимальный режим --- подпериод внутри этапа адаптации со стабильной
работой скважины и лучшими показателями (выделен \BLOCK{ if optimal_regime.source == "auto" }автоматически по composite Score\BLOCK{ else }вручную\BLOCK{ endif }).
Сравнение с этапом наблюдения выполнено именно по этому окну.
}

\vspace{0.3em}

{\scriptsize Табл.~4.3. Параметры оптимального режима.}
\vspace{0.2em}

{\footnotesize
\begin{tabularx}{\textwidth}{|X|R{3.5cm}|}
\hline
\textbf{Параметр} & \textbf{Значение} \\
\hline
Период & \VAR{optimal_regime.date_from}~---~\VAR{optimal_regime.date_to} \\
\hline
Длительность & \VAR{optimal_regime.duration_label} \\
\hline
$P_{\text{тр}}$ (медиана) & \VAR{optimal_regime.p_tube_median_fmt} кгс/см$^2$ \\
\hline
$P_{\text{лин}}$ (медиана) & \VAR{optimal_regime.p_line_median_fmt} кгс/см$^2$ \\
\hline
$\Delta P$ (медиана) & \VAR{optimal_regime.dp_median_fmt} кгс/см$^2$ \\
\hline
$Q$ медианный & \VAR{optimal_regime.flow_median_fmt} тыс.м$^3$/сут \\
\hline
$Q$ средний & \VAR{optimal_regime.flow_avg_fmt} тыс.м$^3$/сут \\
\hline
$Q$ min / max & \VAR{optimal_regime.q_min_fmt} / \VAR{optimal_regime.q_max_fmt} тыс.м$^3$/сут \\
\hline
$\Sigma Q$ накопленный & \VAR{optimal_regime.flow_cumulative_fmt} тыс.м$^3$ \\
\hline
КИВ & \VAR{optimal_regime.utilization_pct_fmt}\% \\
\hline
Продувок & \VAR{optimal_regime.purge_count} \\
\hline
Вбросов реагента & \VAR{optimal_regime.reagent_count} \\
\hline
\BLOCK{ if optimal_regime.source == "auto" and optimal_regime.score is defined }
Score & \VAR{optimal_regime.score_fmt} \\
\hline
\BLOCK{ endif }
\end{tabularx}
}

\BLOCK{ if optimal_regime.narrative }
\vspace{0.4em}
{\footnotesize \VAR{optimal_regime.narrative} }
\BLOCK{ endif }

\BLOCK{ if optimal_regime.combined_chart_path }
\vspace{0.5em}
\begin{center}
\includegraphics[width=\textwidth,height=0.55\textheight,keepaspectratio]{\VAR{optimal_regime.combined_chart_path}}
\end{center}
{\scriptsize Рис.~4.2. Оптимальный режим: давления, $\Delta P$ и Q(t).}
\BLOCK{ endif }

\BLOCK{ endif }

\newpage
\BLOCK{ endif }
%# /4. Этап адаптации

% ═══════════════════════════════════════════════════════════════
%  5. ГРАФИКИ СРАВНЕНИЯ (toggle: include_sections.charts_compare)
% ═══════════════════════════════════════════════════════════════
\BLOCK{ if include_sections.charts_compare }
\section*{5. Графики этапов (сравнение)}

\BLOCK{ if observation.pressure_chart_path or adaptation.pressure_chart_path }
\noindent
\begin{minipage}[t]{0.49\textwidth}
\centering
\BLOCK{ if observation.pressure_chart_path }
\includegraphics[width=\textwidth]{\VAR{observation.pressure_chart_path}}\\[-0.3em]
{\scriptsize Наблюдение}
\BLOCK{ else }
{\scriptsize Нет данных}
\BLOCK{ endif }
\end{minipage}\hfill
\begin{minipage}[t]{0.49\textwidth}
\centering
\BLOCK{ if adaptation.pressure_chart_path }
\includegraphics[width=\textwidth]{\VAR{adaptation.pressure_chart_path}}\\[-0.3em]
{\scriptsize Адаптация}
\BLOCK{ else }
{\scriptsize Нет данных}
\BLOCK{ endif }
\end{minipage}

\vspace{0.3em}
{\scriptsize\centering Рис.~5.1. Давления P\,трубное / P\,линейное. Маркеры: пунктир --- продувки, точки --- вбросы реагента.\par}
\vspace{0.6em}
\BLOCK{ endif }

\BLOCK{ if observation.dp_chart_path or adaptation.dp_chart_path }
\noindent
\begin{minipage}[t]{0.49\textwidth}
\centering
\BLOCK{ if observation.dp_chart_path }
\includegraphics[width=\textwidth]{\VAR{observation.dp_chart_path}}\\[-0.3em]
{\scriptsize Наблюдение}
\BLOCK{ else }
{\scriptsize Нет данных}
\BLOCK{ endif }
\end{minipage}\hfill
\begin{minipage}[t]{0.49\textwidth}
\centering
\BLOCK{ if adaptation.dp_chart_path }
\includegraphics[width=\textwidth]{\VAR{adaptation.dp_chart_path}}\\[-0.3em]
{\scriptsize Адаптация}
\BLOCK{ else }
{\scriptsize Нет данных}
\BLOCK{ endif }
\end{minipage}

\vspace{0.3em}
{\scriptsize\centering Рис.~5.2. Перепад $\Delta P$ с линией тренда.\par}
\vspace{0.6em}
\BLOCK{ endif }

\BLOCK{ if observation.flow_chart_path or adaptation.flow_chart_path }
\noindent
\begin{minipage}[t]{0.49\textwidth}
\centering
\BLOCK{ if observation.flow_chart_path }
\includegraphics[width=\textwidth]{\VAR{observation.flow_chart_path}}\\[-0.3em]
{\scriptsize Наблюдение}
\BLOCK{ else }
{\scriptsize Нет данных}
\BLOCK{ endif }
\end{minipage}\hfill
\begin{minipage}[t]{0.49\textwidth}
\centering
\BLOCK{ if adaptation.flow_chart_path }
\includegraphics[width=\textwidth]{\VAR{adaptation.flow_chart_path}}\\[-0.3em]
{\scriptsize Адаптация}
\BLOCK{ else }
{\scriptsize Нет данных}
\BLOCK{ endif }
\end{minipage}

\vspace{0.3em}
{\scriptsize\centering Рис.~5.3. Мгновенный дебит Q(t), тыс.\,м$^3$/сут.\par}
\BLOCK{ endif }

\newpage
\BLOCK{ endif }
%# /5. Графики сравнения

% ═══════════════════════════════════════════════════════════════
%  6. СРАВНИТЕЛЬНЫЙ АНАЛИЗ  (toggle: include_sections.comparison)
% ═══════════════════════════════════════════════════════════════
\BLOCK{ if include_sections.comparison }
\section*{6. Сравнительный анализ этапов}

\subsection*{6.1. Сравнительная таблица показателей}

{\scriptsize Табл.~6.1. Сравнение основных показателей\BLOCK{ if optimal_regime } (+ колонки оптимального режима)\BLOCK{ endif }.}
\vspace{0.2em}

{\scriptsize
\BLOCK{ if optimal_regime }
\begin{tabularx}{\textwidth}{|P{3.2cm}|R{1.2cm}|R{1.2cm}|R{1.1cm}|R{1.1cm}|R{1.2cm}|R{1.1cm}|R{1.1cm}|}
\hline
\textbf{Показатель} & \textbf{Набл.} & \textbf{Адапт.} & \textbf{$\Delta$} & \textbf{$\Delta$\%} & \textbf{Оптим.} & \textbf{$\Delta$} & \textbf{$\Delta$\%} \\
\hline
$P_{\text{тр}}$, кгс/см$^2$ & \VAR{observation.p_tube_median_fmt} & \VAR{adaptation.p_tube_median_fmt} & \VAR{comparison.delta_p_tube_fmt} & \VAR{comparison.delta_p_tube_pct_fmt} & \VAR{optimal_regime.p_tube_median_fmt} & \VAR{comparison_optimal.delta_p_tube_fmt} & \VAR{comparison_optimal.delta_p_tube_pct_fmt} \\
\hline
$P_{\text{лин}}$, кгс/см$^2$ & \VAR{observation.p_line_median_fmt} & \VAR{adaptation.p_line_median_fmt} & \VAR{comparison.delta_p_line_fmt} & \VAR{comparison.delta_p_line_pct_fmt} & \VAR{optimal_regime.p_line_median_fmt} & \VAR{comparison_optimal.delta_p_line_fmt} & \VAR{comparison_optimal.delta_p_line_pct_fmt} \\
\hline
$\Delta P$, кгс/см$^2$ & \VAR{observation.dp_median_fmt} & \VAR{adaptation.dp_median_fmt} & \VAR{comparison.delta_dp_fmt} & \VAR{comparison.delta_dp_pct_fmt} & \VAR{optimal_regime.dp_median_fmt} & \VAR{comparison_optimal.delta_dp_fmt} & \VAR{comparison_optimal.delta_dp_pct_fmt} \\
\hline
$Q$ мед., тыс.м$^3$/сут & \VAR{observation.flow_median_fmt} & \VAR{adaptation.flow_median_fmt} & \VAR{comparison.delta_flow_median_fmt} & \VAR{comparison.delta_flow_median_pct_fmt} & \VAR{optimal_regime.flow_median_fmt} & \VAR{comparison_optimal.delta_flow_median_fmt} & \VAR{comparison_optimal.delta_flow_median_pct_fmt} \\
\hline
$Q$ ср., тыс.м$^3$/сут & \VAR{observation.flow_avg_fmt} & \VAR{adaptation.flow_avg_fmt} & \VAR{comparison.delta_flow_avg_fmt} & \VAR{comparison.delta_flow_avg_pct_fmt} & \VAR{optimal_regime.flow_avg_fmt} & \VAR{comparison_optimal.delta_flow_avg_fmt} & \VAR{comparison_optimal.delta_flow_avg_pct_fmt} \\
\hline
$Q$ min, тыс.м$^3$/сут & \VAR{observation.q_min_fmt} & \VAR{adaptation.q_min_fmt} & --- & --- & \VAR{optimal_regime.q_min_fmt} & --- & --- \\
\hline
$Q$ max, тыс.м$^3$/сут & \VAR{observation.q_max_fmt} & \VAR{adaptation.q_max_fmt} & --- & --- & \VAR{optimal_regime.q_max_fmt} & --- & --- \\
\hline
КИВ, \% & \VAR{observation.utilization_pct_fmt} & \VAR{adaptation.utilization_pct_fmt} & \VAR{comparison.delta_utilization_pct_fmt} & --- & \VAR{optimal_regime.utilization_pct_fmt} & \VAR{comparison_optimal.delta_utilization_pct_fmt} & --- \\
\hline
Продувок & \VAR{observation.purge_count} & \VAR{adaptation.purge_count} & \VAR{comparison.delta_purge_count_fmt} & --- & \VAR{optimal_regime.purge_count} & --- & --- \\
\hline
Длительность & \VAR{observation.duration_label} & \VAR{adaptation.duration_label} & --- & --- & \VAR{optimal_regime.duration_label} & --- & --- \\
\hline
\end{tabularx}
\BLOCK{ else }
\begin{tabularx}{\textwidth}{|X|R{2cm}|R{2cm}|R{2cm}|R{2cm}|}
\hline
\textbf{Показатель} & \textbf{Наблюдение} & \textbf{Адаптация} & \textbf{$\Delta$} & \textbf{$\Delta$, \%} \\
\hline
$P_{\text{тр}}$, кгс/см$^2$ & \VAR{observation.p_tube_median_fmt} & \VAR{adaptation.p_tube_median_fmt} & \VAR{comparison.delta_p_tube_fmt} & \VAR{comparison.delta_p_tube_pct_fmt} \\
\hline
$P_{\text{лин}}$, кгс/см$^2$ & \VAR{observation.p_line_median_fmt} & \VAR{adaptation.p_line_median_fmt} & \VAR{comparison.delta_p_line_fmt} & \VAR{comparison.delta_p_line_pct_fmt} \\
\hline
$\Delta P$, кгс/см$^2$ & \VAR{observation.dp_median_fmt} & \VAR{adaptation.dp_median_fmt} & \VAR{comparison.delta_dp_fmt} & \VAR{comparison.delta_dp_pct_fmt} \\
\hline
$Q$ мед., тыс.м$^3$/сут & \VAR{observation.flow_median_fmt} & \VAR{adaptation.flow_median_fmt} & \VAR{comparison.delta_flow_median_fmt} & \VAR{comparison.delta_flow_median_pct_fmt} \\
\hline
$Q$ ср., тыс.м$^3$/сут & \VAR{observation.flow_avg_fmt} & \VAR{adaptation.flow_avg_fmt} & \VAR{comparison.delta_flow_avg_fmt} & \VAR{comparison.delta_flow_avg_pct_fmt} \\
\hline
КИВ, \% & \VAR{observation.utilization_pct_fmt} & \VAR{adaptation.utilization_pct_fmt} & \VAR{comparison.delta_utilization_pct_fmt} & --- \\
\hline
Продувок & \VAR{observation.purge_count} & \VAR{adaptation.purge_count} & \VAR{comparison.delta_purge_count_fmt} & --- \\
\hline
Длительность & \VAR{observation.duration_label} & \VAR{adaptation.duration_label} & --- & --- \\
\hline
\end{tabularx}
\BLOCK{ endif }
}

\vspace{0.5em}

% ── 6.1.1 — Сравнение адаптации с зафиксированными baseline-ами ──
\BLOCK{ if adapt_vs_baseline and (adapt_vs_baseline.baselines.customer or adapt_vs_baseline.baselines.observation) }
\subsection*{6.1.1. Сравнение этапа адаптации с baseline-ами}

{\footnotesize
Сравнение медианных показателей этапа адаптации с двумя зафиксированными
опорными точками: <<baseline заказчика>> (по сводке УзКорГаз, well\_daily)
и <<baseline наблюдения>> (НАШИ данные за этап наблюдения, pressure\_raw +
flow\_rate). Если соответствующий baseline не зафиксирован — прочерк.
}

\vspace{0.2em}

{\scriptsize
\begin{tabularx}{\textwidth}{|P{3.4cm}|R{1.6cm}|R{1.6cm}|R{1.5cm}|R{1.6cm}|R{1.5cm}|R{1.6cm}|}
\hline
\textbf{Показатель} & \textbf{Адапт.}
& \textbf{Base заказч.} & \textbf{$\Delta$ vs заказч.}
& \textbf{Base набл.}   & \textbf{$\Delta$ vs набл.}
& \textbf{Прим.} \\
\hline
%% set bc = adapt_vs_baseline.baselines.customer
%% set bo = adapt_vs_baseline.baselines.observation
%% set dc = adapt_vs_baseline.customer
%% set do = adapt_vs_baseline.observation

$Q$ мед., тыс.м$^3$/сут & \VAR{adaptation.flow_median_fmt}
& \BLOCK{ if bc and bc.q_total_median is not none }\VAR{'%.2f'|format(bc.q_total_median)}\BLOCK{ else }---\BLOCK{ endif }
& \BLOCK{ if dc.delta_q_med is not none }\VAR{'%+.2f'|format(dc.delta_q_med)}\BLOCK{ else }---\BLOCK{ endif }
& \BLOCK{ if bo and bo.q_total_median is not none }\VAR{'%.2f'|format(bo.q_total_median)}\BLOCK{ else }---\BLOCK{ endif }
& \BLOCK{ if do.delta_q_med is not none }\VAR{'%+.2f'|format(do.delta_q_med)}\BLOCK{ else }---\BLOCK{ endif }
& \BLOCK{ if dc.delta_q_med_pct is not none }\VAR{'%+.0f'|format(dc.delta_q_med_pct)}\%\BLOCK{ endif } \\
\hline
$\Delta P$ мед., кгс/см$^2$ & \VAR{adaptation.dp_median_fmt}
& \BLOCK{ if bc and bc.dp_median is not none }\VAR{'%.2f'|format(bc.dp_median)}\BLOCK{ else }---\BLOCK{ endif }
& \BLOCK{ if dc.delta_dp_med is not none }\VAR{'%+.2f'|format(dc.delta_dp_med)}\BLOCK{ else }---\BLOCK{ endif }
& \BLOCK{ if bo and bo.dp_median is not none }\VAR{'%.2f'|format(bo.dp_median)}\BLOCK{ else }---\BLOCK{ endif }
& \BLOCK{ if do.delta_dp_med is not none }\VAR{'%+.2f'|format(do.delta_dp_med)}\BLOCK{ else }---\BLOCK{ endif }
& \BLOCK{ if dc.delta_dp_med_pct is not none }\VAR{'%+.0f'|format(dc.delta_dp_med_pct)}\%\BLOCK{ endif } \\
\hline
$P_{\text{тр}}$ мед., кгс/см$^2$ & \VAR{adaptation.p_tube_median_fmt}
& \BLOCK{ if bc and bc.p_wellhead_median is not none }\VAR{'%.2f'|format(bc.p_wellhead_median)}\BLOCK{ else }---\BLOCK{ endif }
& \BLOCK{ if dc.delta_p_tube_med is not none }\VAR{'%+.2f'|format(dc.delta_p_tube_med)}\BLOCK{ else }---\BLOCK{ endif }
& \BLOCK{ if bo and bo.p_wellhead_median is not none }\VAR{'%.2f'|format(bo.p_wellhead_median)}\BLOCK{ else }---\BLOCK{ endif }
& \BLOCK{ if do.delta_p_tube_med is not none }\VAR{'%+.2f'|format(do.delta_p_tube_med)}\BLOCK{ else }---\BLOCK{ endif }
& --- \\
\hline
$P_{\text{лин}}$ мед., кгс/см$^2$ & \VAR{adaptation.p_line_median_fmt}
& \BLOCK{ if bc and bc.p_flowline_median is not none }\VAR{'%.2f'|format(bc.p_flowline_median)}\BLOCK{ else }---\BLOCK{ endif }
& \BLOCK{ if dc.delta_p_line_med is not none }\VAR{'%+.2f'|format(dc.delta_p_line_med)}\BLOCK{ else }---\BLOCK{ endif }
& \BLOCK{ if bo and bo.p_flowline_median is not none }\VAR{'%.2f'|format(bo.p_flowline_median)}\BLOCK{ else }---\BLOCK{ endif }
& \BLOCK{ if do.delta_p_line_med is not none }\VAR{'%+.2f'|format(do.delta_p_line_med)}\BLOCK{ else }---\BLOCK{ endif }
& --- \\
\hline
\end{tabularx}
}

\vspace{0.3em}
{\scriptsize
\textit{
\BLOCK{ if bc }Baseline заказч.: <<\VAR{bc.name}>>, \VAR{bc.period_from}--\VAR{bc.period_to}.\BLOCK{ endif }
\BLOCK{ if bo }Baseline набл.: <<\VAR{bo.name}>>, \VAR{bo.period_from}--\VAR{bo.period_to}.\BLOCK{ endif }
}
}

\vspace{0.5em}
\BLOCK{ endif }

\subsection*{6.2. Выводы}

\BLOCK{ if conclusions }
{\footnotesize
\begin{enumerate}\setlength{\itemsep}{-1pt}
%% for c in conclusions
\item \VAR{c}
%% endfor
\end{enumerate}
}
\BLOCK{ else }
{\footnotesize Недостаточно данных для формирования выводов.}
\BLOCK{ endif }

\vspace{1em}
\BLOCK{ endif }
%# /6. Сравнительный анализ

% ═══════════════════════════════════════════════════════════════
%  ПОДПИСИ
% ═══════════════════════════════════════════════════════════════
\vfill

{\footnotesize
\begin{tabularx}{\textwidth}{XX}
\textbf{Исполнитель} & \textbf{Заказчик} \\[2em]
\rule{5cm}{0.4pt} & \rule{5cm}{0.4pt} \\
(подпись, расшифровка) & (подпись, расшифровка) \\[1em]
\rule{5cm}{0.4pt} & \rule{5cm}{0.4pt} \\
(дата) & (дата) \\
\end{tabularx}
}

\end{document}
